% !PRINTED 479
\weberpage{479}

Hereupon, if $\varrho$ is expressed in terms of the roots $\alpha$, we can by hypothesis apply the permutation $\pi$, and thus, if $\varrho$ thereby passes into $\varrho_a$, we obtain $g(\varrho_a) = 0$; that is, $\varrho_a$ is also a root of the resolvent, and the permutation $\pi$ therefore corresponds to a substitution $(\varrho, \varrho_a)$ of the field $\Omega(\varrho)$, i.\,e. $\pi$ belongs to the \gesperrt{Galois group}.

From this we conclude further the following important theorem:

d) \gesperrt{If $P$ is a group of permutations of the $m$ roots $\alpha$ which possesses the properties a) and b), then $P$ is the Galois group of the equation $F(x) = 0$.}

For, in the first place, by c) every permutation of $P$ belongs to the \gesperrt{Galois group}, and $P$ is therefore certainly a divisor of the latter.

But if $P$ comprises only $\nu$ permutations, we use the following notation.
\weq{5}{\pi_1, \pi_2 \ldots \pi_\nu}
If we apply these permutations to $\varrho$, let the result be:
\weq{6}{\varrho_1, \varrho_2 \ldots \varrho_\nu.}
If we apply one of the permutations (5), say $\pi_k$, to one of the quantities (6), say to $\varrho_i$, the effect is the same as if $\pi_i \pi_k$ had been applied to $\varrho$; let the result of this permutation be $\varrho_i'$. Now it lies in the hypothesis that $P$ is a group that the composed permutation $\pi_i \pi_k$ also belongs to $P$, and hence that $\varrho_i'$ is contained among the quantities (6). At the same time, by \S.\,146, 2., $\varrho_i'$ is different from $\varrho_h'$ as soon as $\varrho_i$ is different from $\varrho_h$. From this it follows that the quantities
\weq{7}{\varrho_1', \varrho_2' \ldots \varrho_\nu'}
agree with the quantities (6), apart from the order.

The product
\[
g'(t) = (t - \varrho_1)(t - \varrho_2) \ldots (t - \varrho_\nu)
\]
therefore remains unchanged by the permutations of the group $P$, and is consequently, since we presuppose the property b) of $P$, a rational function of $t$ in $\Omega$. At the same time $g'(t)$ is a divisor of $g(t)$, and must therefore, since $g(t)$ is irreducible, agree with $g(t)$; hence $\nu$ is not smaller than the degree of the \gesperrt{Galois group}, and $P$ is identical with the \gesperrt{Galois group}.
